% Theorem 2 (Entropy-Solution Stability -- Baseline Rate) -- Appendix Proof
% Included under \subsection{Proof of Theorem~\ref{thm:stability}} in A1_proofs.tex

\begin{theorem}[Entropy-solution stability --- restated]
\label{thm:stability-app}
Let $\physsol$ be the Kruzhkov entropy solution of the scalar conservation law $\dt \mathbf{u} + \dx \flux(\mathbf{u}) = 0$ at a fixed physical time $T_{\mathrm{phys}}$. View the solution as a state vector $u \in \R^{d}$ ($d = N_{x}$ is the number of spatial grid points) and define the target distribution
\begin{equation}
\rhotrue \coloneqq \Law(\physsol).
\label{eq:thm2-rho-star}
\end{equation}

The forward diffusion process (VE-SDE, schedule $\sigtau$):
\begin{equation}
\ptau \coloneqq \rhotrue \ast \Gtau, \qquad
\Gtau(u) = (4\pi\tau)^{-d/2} \exp\!\left(-\frac{\norm{u}^{2}}{4\tau}\right), \quad \tau \in [0, T_{d}].
\label{eq:thm2-ptau}
\end{equation}

The true score field:
\begin{equation}
\score(u, \tau) \coloneqq \nabla_{u} \log \ptau(u).
\label{eq:thm2-score}
\end{equation}

The reverse ODE (standard EDM form, $\tau$ integrated from $T_{d}$ to $0$):

True trajectory:
\begin{align}
\frac{d\unat}{d\tau} &= b(\unat, \tau), \qquad
b(u, \tau) \coloneqq -\sigtau\,\dot\sigma(\tau)\,\score(u, \tau),
\quad \unat(T_{d}) \sim p_{T_{d}},
\label{eq:thm2-reverse-true}
\end{align}

Network-approximated trajectory:
\begin{align}
\frac{d\uhat}{d\tau} &= \hat b(\uhat, \tau), \qquad
\hat b(u, \tau) \coloneqq -\sigtau\,\dot\sigma(\tau)\,\sth(u, \tau),
\quad \uhat(T_{d}) \sim p_{T_{d}}.
\label{eq:thm2-reverse-approx}
\end{align}

Notation: $\unat(\tau)$ is the true-score-driven reverse trajectory; $\uhat(\tau)$ is the learned-score-driven reverse trajectory. $\norm{\cdot}$ denotes the Euclidean norm on $\R^{d}$. $\Wass{1}(\mu, \nu)$ is the 1-Wasserstein distance (Kantorovich--Rubinstein dual). The expectation $\E$ is taken over randomness (initial noise $+$ learned score approximation).

\medskip
\noindent
\textbf{Main theorem.} Under Assumptions~(A1)--(A5), the generated sample distribution $\Law(\uhat(0))$ and the target distribution $\rhotrue$ satisfy
\begin{equation}
\boxed{\;
\Wass{1}\!\bigl(\Law(\uhat(0)),\; \rhotrue\bigr)
\;\le\; C\,\varepsilon\,\exp(\amp T_{d})
\;},
\label{eq:thm2-main}
\end{equation}
where $\amp \coloneqq \sup_{\tau \in [0, T_{d}]} \abs{\sigtau\,\dot\sigma(\tau)}\,L_{s}(\tau)$, $L_{s}(\tau)$ is the Lipschitz constant of $\nabla_{u} \score(\cdot, \tau)$, and $\varepsilon$ is the score-matching training error bound. The constant $C = C_{\sigma} \sqrt{T_{d}}\,(1 + C_{\mathrm{stab}})$ depends only on the schedule upper bound $C_{\sigma}$ and the distributional stability constant $C_{\mathrm{stab}}$, and is independent of $\varepsilon$ and $\amp$.

\medskip
\noindent
\textbf{Connection to the PDE solution.} Equation~\eqref{eq:thm2-main} provides a quantitative guarantee that the diffusion model's sampled solution distribution is close to the PDE solution distribution. If the training data originates from a numerical solver, the Kruzhkov $L^{1}$-contractivity ensures that the numerical error does not amplify during physical time propagation; the bound~\eqref{eq:thm2-main} transfers stably to the true PDE solution (see \S\,PDE Connection below).
\end{theorem}

\paragraph{Assumptions.}
\begin{enumerate}
\item[(A1)] \textbf{Score gradient Lipschitz regularity.} For each $\tau > 0$, the true score field $\score(\cdot, \tau)$ is $C^{1}$ and its Jacobian's operator norm is uniformly bounded: there exists $L_{s}(\tau) < \infty$ such that
\begin{equation}
\norm{\nabla_{u} \score(u, \tau)}_{\mathrm{op}} \le L_{s}(\tau), \quad \forall\, u \in \R^{d},\; \tau \in (0, T_{d}].
\label{eq:thm2-assum-A1}
\end{equation}
Define $\amp(\tau) \coloneqq \abs{\sigtau\,\dot\sigma(\tau)}\,L_{s}(\tau)$ and set $\amp \coloneqq \sup_{\tau \in [0, T_{d}]} \amp(\tau) < \infty$.

\textit{Verifiability:} Since $\ptau = \rhotrue \ast \Gtau$ and $\Gtau$ is the heat kernel, for any $\tau > 0$, $\ptau \in C^{\infty}$ and $\nabla_{u} \score = \nabla_{u}^{2} \log \ptau$ is uniformly bounded. $\rhotrue$ is a pushforward of a BV measure (containing Dirac point masses at shocks); smoothness holds after convolution. $L_{s}(\tau) \sim 1/\tau$ may diverge as $\tau \to 0$, but we only require $\sup_{\tau \in [\tau_{\mathrm{end}}, T_{d}]} L_{s}(\tau) < \infty$ (early-stopping $\tau_{\mathrm{end}} > 0$), consistent with \citet{sarkar2026scoreshocks}.

\item[(A2)] \textbf{Noise schedule boundedness.} There exists $C_{\sigma} > 0$ such that
\begin{equation}
\abs{\sigtau\,\dot\sigma(\tau)} \le C_{\sigma}, \quad \forall\, \tau \in [0, T_{d}].
\label{eq:thm2-assum-A2}
\end{equation}

\textit{Verifiability:} For VE-type schedules (e.g., $\sigma^{2}(\tau) = 2\physvis \tau$), $\sigma \dot\sigma = \physvis$ is constant, which directly satisfies this.

\item[(A3)] \textbf{Score training error (weighted $L^{2}$ bound).} The learned score network $\sth$ satisfies, under the true data distribution $\ptau$,
\begin{equation}
\int_{0}^{T_{d}} \E_{u \sim \ptau}\!\bigl[\,\norm{\sth(u, \tau) - \score(u, \tau)}^{2}\,\bigr] \, d\tau \;\le\; \varepsilon^{2}.
\label{eq:thm2-assum-A3}
\end{equation}

\textit{Verifiability:} This is the training objective of denoising score matching \citep{vincent2011connection} (with weight $\sigtau^{2}$); $\varepsilon$ is determined by the training loss.

\item[(A4)] \textbf{Distributional stability (Kantorovich--Rubinstein control).} There exists a constant $C_{\mathrm{stab}} > 0$ such that for any $\tau \in [0, T_{d}]$ and any Lipschitz function $f: \R^{d} \to \R$,
\begin{equation}
\abs{\,\E_{\uhat(\tau)}[f(\uhat)] - \E_{\unat(\tau)}[f(\unat)]\,}
\;\le\;
C_{\mathrm{stab}} \cdot \Wass{1}\!\bigl(\Law(\uhat(\tau)),\, \Law(\unat(\tau))\bigr) \cdot \Lip(f),
\label{eq:thm2-assum-A4}
\end{equation}
where $\Lip(f) \coloneqq \sup_{u \neq v} \abs{f(u) - f(v)} / \norm{u - v}$.

\textit{Verifiability:} This is a direct generalization of the Kantorovich--Rubinstein dual $\abs{\int f d\mu - \int f d\nu} \le \Lip(f) \cdot \Wass{1}(\mu, \nu)$, for which $C_{\mathrm{stab}} = 1$. Introducing $C_{\mathrm{stab}}$ accommodates additional smoothness offsets introduced by the score network; under compactly supported measures and smooth score assumptions, $C_{\mathrm{stab}} = \mathcal{O}(1)$.

\item[(A5)] \textbf{Network Lipschitz bound.} The learned score network $\sth(\cdot, \tau)$ is uniformly Lipschitz across all $\tau \in [0, T_{d}]$: there exists $L_{s_{\theta}} < \infty$ independent of $\tau$ such that for all $u, v \in \R^{d}$,
\begin{equation}
\norm{\sth(u, \tau) - \sth(v, \tau)} \le L_{s_{\theta}} \norm{u - v}.
\label{eq:thm2-assum-A5}
\end{equation}

\textit{Verifiability:} Standard MLPs/UNets are Lipschitz on compact input domains; boundedness of weights is ensured by spectral normalization or training regularization.
\end{enumerate}

\begin{proof}
The proof proceeds in nine steps. Steps~1--3 set up the synchronous coupling framework; Steps~4--5 address the core technical difficulty (measure mismatch); Steps~6--8 synthesize the Gronwall estimate; Step~9 yields the final $\Wass{1}$ bound.

\medskip
\noindent
\textbf{Step~1 --- Synchronous coupling.}
Define the error process
\begin{equation}
\errsc(\tau) \coloneqq \uhat(\tau) - \unat(\tau), \qquad \tau \in [0, T_{d}].
\label{eq:thm2-step1}
\end{equation}

Since both trajectories share the same initial distribution $p_{T_{d}}$, we have $\errsc(T_{d}) = 0$ almost surely. Differentiating along the reverse ODE with respect to $\tau$ (with $\tau$ decreasing from $T_{d}$ to $0$):
\begin{equation}
\frac{d}{d\tau} \errsc(\tau)
= \hat b(\uhat(\tau), \tau) - b(\unat(\tau), \tau).
\label{eq:thm2-step1-ode}
\end{equation}

Add-and-subtract decomposition:
\begin{align}
\frac{d}{d\tau} \errsc(\tau)
&= \bigl[b(\uhat(\tau), \tau) - b(\unat(\tau), \tau)\bigr]
   \;+\;
   \bigl[\hat b(\uhat(\tau), \tau) - b(\uhat(\tau), \tau)\bigr] \notag \\[4pt]
&\eqqcolon (\mathrm{I}) + (\mathrm{II}).
\label{eq:thm2-step1-decomp}
\end{align}

The two terms differ in nature:
\begin{itemize}
\item $(\mathrm{I})$: propagation error due to trajectory deviation $\uhat \neq \unat$ under the \emph{same} drift field $b$;
\item $(\mathrm{II})$: model error due to the score approximation $\sth \neq \score$ in the drift field itself.
\end{itemize}

\medskip
\noindent
\textbf{Step~2 --- Controlling $(\mathrm{I})$ (propagation error).}
From $b(u, \tau) = -\sigtau\,\dot\sigma(\tau)\,\score(u, \tau)$ and Assumption~(A1):
\begin{align}
\norm{b(\uhat, \tau) - b(\unat, \tau)}
&\le \abs{\sigtau\,\dot\sigma(\tau)} \cdot \norm{\score(\uhat, \tau) - \score(\unat, \tau)} \notag \\
&\le \abs{\sigtau\,\dot\sigma(\tau)} \cdot L_{s}(\tau) \cdot \norm{\uhat - \unat} \notag \\
&= \amp(\tau)\,\norm{\errsc(\tau)}
\le \amp\,\norm{\errsc(\tau)}.
\label{eq:thm2-step2}
\end{align}

Here $\amp = \sup_{\tau \in [0, T_{d}]} \amp(\tau)$ is precisely the Score Shocks amplification exponent defined in \citet{sarkar2026scoreshocks}. This is the \textbf{origin} of the $\exp(\amp T_{d})$ factor --- the Lipschitz constant of the drift field $b$ equals $\amp$, so errors propagate at exponential rate $\amp$.

\medskip
\noindent
\textbf{Step~3 --- Controlling $(\mathrm{II})$ (score model error).}
From the definitions of $b, \hat b$ and Assumption~(A2):
\begin{align}
\norm{\hat b(\uhat, \tau) - b(\uhat, \tau)}
&= \abs{\sigtau\,\dot\sigma(\tau)} \cdot \norm{\sth(\uhat, \tau) - \score(\uhat, \tau)} \notag \\
&\le C_{\sigma}\,\norm{\sth(\uhat, \tau) - \score(\uhat, \tau)}.
\label{eq:thm2-step3}
\end{align}

\medskip
\noindent
\textbf{Step~4 --- Expectation inequality.}
Take the Euclidean norm and expectation of~\eqref{eq:thm2-step1-decomp}, and define
\begin{equation}
\delta(\tau) \coloneqq \E\,\norm{\errsc(\tau)}.
\label{eq:thm2-step4-delta}
\end{equation}

By Jensen's inequality $\E\norm{\frac{d}{d\tau}\errsc} \ge \abs{\frac{d}{d\tau}\E\norm{\errsc}}$ (strictly valid under absolute continuity; handled via the chain rule; see Remark~1), combining~\eqref{eq:thm2-step2} and~\eqref{eq:thm2-step3}:
\begin{equation}
\delta'(\tau)
\;\le\; \amp\,\delta(\tau)
\;+\; C_{\sigma}\;\E_{\uhat(\tau)}\!\bigl[\norm{\sth(\uhat, \tau) - \score(\uhat, \tau)}\bigr],
\label{eq:thm2-step4-ineq}
\end{equation}
where $\delta'(\tau) = d\delta/d\tau$ along the reverse time direction ($\tau \searrow 0$).

\begin{remark}
\textbf{(Direction of Jensen's inequality).} Strictly, $\frac{d}{d\tau}\E\norm{\errsc} = \E\bigl[\frac{\errsc}{\norm{\errsc}} \cdot \frac{d \errsc}{d\tau}\bigr] \le \E\norm{\frac{d \errsc}{d\tau}}$, so the inequality direction is correct. On the null set where $\norm{\errsc(\tau)} = 0$, the inequality degenerates to $0 \le \E\norm{\frac{d \errsc}{d\tau}}$, which does not affect the integral bound.
\end{remark}

\medskip
\noindent
\textbf{Step~5 --- Repairing the measure mismatch (core technical step).}

\noindent
\textbf{Problem:} The second expectation in~\eqref{eq:thm2-step4-ineq} is taken under the distribution of $\uhat(\tau)$ (the learned trajectory), whereas the training error Assumption~(A3) is controlled under $\ptau$ (the distribution of the true score trajectory $\unat(\tau)$). The two \textbf{measures differ}.

\noindent
\textbf{Repair via triangle decomposition.} Define
\begin{equation}
g(u, \tau) \coloneqq \norm{\sth(u, \tau) - \score(u, \tau)}.
\label{eq:thm2-step5-g}
\end{equation}
Then
\begin{align}
\E_{\uhat(\tau)}[g(\uhat, \tau)]
&= \E_{\unat(\tau)}[g(\unat, \tau)]
   \;+\;
   \Bigl(\E_{\uhat(\tau)}[g(\uhat, \tau)] - \E_{\unat(\tau)}[g(\unat, \tau)]\Bigr) \notag \\[4pt]
&\le \E_{\unat(\tau)}[g(\unat, \tau)]
   \;+\;
   \abs{\,\E_{\uhat(\tau)}[g(\uhat, \tau)] - \E_{\unat(\tau)}[g(\unat, \tau)]\,}.
\label{eq:thm2-step5-decomp}
\end{align}

By Assumption~(A4) (distributional stability), applied to the function $g(\cdot, \tau)$:
\begin{equation}
\abs{\,\E_{\uhat(\tau)}[g] - \E_{\unat(\tau)}[g]\,}
\;\le\;
C_{\mathrm{stab}} \cdot \Wass{1}\!\bigl(\Law(\uhat(\tau)),\, \Law(\unat(\tau))\bigr) \cdot \Lip(g).
\label{eq:thm2-step5-A4}
\end{equation}

By construction of the synchronous coupling,
\begin{equation}
\Wass{1}\!\bigl(\Law(\uhat(\tau)), \Law(\unat(\tau))\bigr) \le \E\norm{\uhat(\tau) - \unat(\tau)} = \delta(\tau).
\label{eq:thm2-step5-W1}
\end{equation}

By Assumptions~(A5) and~(A1), the Lipschitz constant of $g(\cdot, \tau) = \norm{\sth(\cdot, \tau) - \score(\cdot, \tau)}$ satisfies
\begin{equation}
\Lip(g(\cdot, \tau)) \;\le\; L_{s_{\theta}} + L_{s}(\tau) \;\le\; L_{s_{\theta}} + \sup_{\tau} L_{s}(\tau) \;\eqqcolon\; L_{g} < \infty,
\label{eq:thm2-step5-Lg}
\end{equation}
where finiteness is guaranteed by early-stopping $\tau_{\mathrm{end}} > 0$ (see the remark following Assumption~A1).

Combining these:
\begin{equation}
\E_{\uhat(\tau)}[g(\uhat, \tau)]
\;\le\;
\E_{\unat(\tau)}[g(\unat, \tau)]
\;+\;
C_{\mathrm{stab}}\,L_{g}\;\delta(\tau).
\label{eq:thm2-step5-final}
\end{equation}

Define $C' \coloneqq C_{\mathrm{stab}}\,L_{g}$.

\medskip
\noindent
\textbf{Step~6 --- Assembling the differential inequality.}
Substituting~\eqref{eq:thm2-step5-final} into~\eqref{eq:thm2-step4-ineq}:
\begin{equation}
\boxed{\;
\delta'(\tau)
\;\le\; (\amp + C')\,\delta(\tau)
\;+\; C_{\sigma}\;\E_{\unat(\tau)}\!\bigl[\norm{\sth(\unat, \tau) - \score(\unat, \tau)}\bigr]
\;}.
\label{eq:thm2-step6}
\end{equation}

This is a standard first-order linear differential inequality (Gronwall type), where the distribution of $\unat(\tau)$ is precisely $\ptau$, matching the measure in the training error Assumption~(A3) after Step~5's repair.

\medskip
\noindent
\textbf{Step~7 --- Score error integral bound (Cauchy--Schwarz).}
Define
\begin{equation}
\beta(\tau) \coloneqq \E_{\unat(\tau)}\!\bigl[\norm{\sth(\unat, \tau) - \score(\unat, \tau)}\bigr].
\label{eq:thm2-step7-beta}
\end{equation}

By Cauchy--Schwarz:
\begin{align}
\int_{0}^{T_{d}} \beta(\tau)\,d\tau
&\le \sqrt{T_{d}}\;\left(\int_{0}^{T_{d}} \E_{\unat(\tau)}\!\bigl[\norm{\sth - \score}^{2}\bigr]\,d\tau\right)^{\!1/2} \notag \\
&\le \sqrt{T_{d}}\;\varepsilon,
\label{eq:thm2-step7-cs}
\end{align}
where the last step uses Assumption~(A3). The key point: the expectation in~(A3) is taken under the \textbf{true trajectory distribution} $\ptau = \Law(\unat(\tau))$, which matches the measure in~\eqref{eq:thm2-step6} after the Step~5 repair.

\medskip
\noindent
\textbf{Step~8 --- Gronwall reverse-time integration.}
Apply the Gronwall argument to the differential inequality~\eqref{eq:thm2-step6} over reverse time $[T_{d}, 0]$. Set $\bar\amp \coloneqq \amp + C'$. Rewrite~\eqref{eq:thm2-step6} as
\begin{equation}
\frac{d}{d\tau}\!\bigl[\delta(\tau)\,e^{-\bar\amp \tau}\bigr]
\;\le\; C_{\sigma}\,\beta(\tau)\,e^{-\bar\amp \tau},
\label{eq:thm2-step8-rewrite}
\end{equation}
where the derivative is along the reverse direction ($\tau$ decreasing). Integrating forward over $[0, T_{d}]$ (equivalently: writing the reverse Gronwall in forward form):
\begin{equation}
\delta(0)\,e^{-\bar\amp \cdot 0} - \delta(T_{d})\,e^{-\bar\amp T_{d}}
\;\le\; C_{\sigma} \int_{0}^{T_{d}} \beta(\tau)\,e^{-\bar\amp \tau}\,d\tau.
\label{eq:thm2-step8-int}
\end{equation}

Since the initial conditions coincide at the noise endpoint ($\uhat(T_{d}) \sim p_{T_{d}}$, $\unat(T_{d}) \sim p_{T_{d}}$), we have $\delta(T_{d}) = 0$. Hence
\begin{align}
\delta(0)
&\le C_{\sigma}\,e^{\bar\amp T_{d}} \int_{0}^{T_{d}} \beta(\tau)\,e^{-\bar\amp \tau}\,d\tau \notag \\
&\le C_{\sigma}\,e^{\bar\amp T_{d}} \int_{0}^{T_{d}} \beta(\tau)\,d\tau.
\label{eq:thm2-step8-delta0}
\end{align}

Substituting~\eqref{eq:thm2-step7-cs}:
\begin{equation}
\delta(0) \;\le\; C_{\sigma} \sqrt{T_{d}}\;\varepsilon\;e^{(\amp + C')T_{d}}.
\label{eq:thm2-step8-final}
\end{equation}

\medskip
\noindent
\textbf{Step~9 --- $\Wass{1}$ upper bound and absorption of constants.}
By the synchronous coupling construction, the $\Wass{1}$ distance is bounded above by the $L^{1}$ expectation under the coupling:
\begin{equation}
\Wass{1}\!\bigl(\Law(\uhat(0)),\, \rhotrue\bigr)
\;\le\; \E\norm{\uhat(0) - \unat(0)}
\;=\; \delta(0).
\label{eq:thm2-step9-W1}
\end{equation}

Note: strictly, $\Wass{1}$ above couples $\Law(\uhat(0))$ vs.\ $\Law(\unat(0))$, and $\Law(\unat(0)) = \rhotrue$ (since the true reverse ODE exactly recovers the target distribution). Therefore $\Wass{1}(\Law(\uhat(0)), \rhotrue) \le \delta(0)$ holds.

Substituting~\eqref{eq:thm2-step8-final} and setting $C \coloneqq C_{\sigma} \sqrt{T_{d}}\,e^{C' T_{d}}$:
\begin{equation}
\Wass{1}\!\bigl(\Law(\uhat(0)),\, \rhotrue\bigr)
\;\le\; C\,\varepsilon\,e^{\amp T_{d}}.
\label{eq:thm2-step9-final}
\end{equation}

Observe that $C' = C_{\mathrm{stab}}\,L_{g}$, where $L_{g}$ depends on the schedule and network structure but is independent of $\varepsilon$. In the asymptotic regime $\varepsilon \ll 1$, $e^{C'T_{d}} = \mathcal{O}(1)$ and can be absorbed into the constant $C$. The final form is as stated in~\eqref{eq:thm2-main}.
\end{proof}

\medskip
\noindent
\textbf{PDE Connection --- Stable transfer from numerical solution to true PDE solution.}
The above proof is carried out entirely within the diffusion model domain (state space $\R^{d}$) and yields a $\Wass{1}$ bound for $\Law(\uhat(0))$ approximating $\rhotrue = \Law(\physsol)$. The following ensures that this bound \textbf{stably transfers} to the true PDE solution when the training data originates from a numerical solver, without amplification during physical time propagation.

\begin{proposition}[Kruzhkov $L^{1}$-contractivity]
Let $\mathbf{u}, \mathbf{v}$ be two Kruzhkov entropy solutions of the same scalar conservation law $\dt \mathbf{u} + \dx \flux(\mathbf{u}) = 0$ with initial data $\mathbf{u}_{0}, \mathbf{v}_{0}$ respectively. Then at any physical time $t > 0$,
\begin{equation}
\norm{\mathbf{u}(\cdot, t) - \mathbf{v}(\cdot, t)}_{L^{1}(\Omega)}
\;\le\;
\norm{\mathbf{u}_{0} - \mathbf{v}_{0}}_{L^{1}(\Omega)}.
\label{eq:thm2-pde-L1-contraction}
\end{equation}
(See \citet{kruzhkov1970first}, Theorem~1, or \citet{dafermos2016hyperbolic}, \S\,6.3.)
\end{proposition}

\begin{corollary}[Numerical error does not amplify]
If the training data is generated by a numerical solution $\mathbf{u}_{h}$ satisfying $\norm{\mathbf{u}_{h} - \physsol}_{L^{1}} \le \delta_{\mathrm{num}}$, then under the discrete correspondence $\norm{u} \approx \sqrt{\Delta x}\,\norm{\mathbf{u}}_{L^{2}}$,
\begin{equation}
\norm{\uhat(0) - u_{h}}
\;\le\;
\underbrace{\norm{\uhat(0) - u^{\star}}}_{\text{Theorem~2 bound}\; \mathcal{O}(\varepsilon e^{\amp T_{d}})}
\;+\;
\underbrace{\norm{u^{\star} - u_{h}}}_{\text{numerical error}\; \mathcal{O}(\delta_{\mathrm{num}})}.
\label{eq:thm2-pde-transfer}
\end{equation}
By~\eqref{eq:thm2-pde-L1-contraction}, $\delta_{\mathrm{num}}$ does not grow with physical time $T_{\mathrm{phys}}$ --- this is a core property distinguishing conservation law entropy solutions from general PDEs (e.g., Navier--Stokes $L^{2}$ errors may propagate in time, whereas conservation law $L^{1}$ errors are bounded by the initial data).

\textit{Practical implication:} In experiments, high-resolution Godunov/WENO solvers are used to generate training reference solutions; $\delta_{\mathrm{num}}$ is much smaller than the score-matching error $\varepsilon$ (typical grid resolution $N_{x} \ge 2048$ yields $\delta_{\mathrm{num}} \sim 10^{-4}$), so the Theorem~2 bound is the dominant error term.
\end{corollary}

\medskip
\noindent
\textbf{Remarks on the form and limitations of the bound.}

\begin{remark}[Origin and tightness of $\exp(\amp)$]
\label{rem:exp-origin}
The exponential factor $\exp(\amp T_{d})$ in~\eqref{eq:thm2-main} originates from the Lipschitz constant $\amp = \sup_{\tau} \abs{\sigma \dot\sigma}\,L_{s}(\tau)$ of the drift field $b$. When the target distribution contains shocks (i.e., $\rhotrue$ has Dirac point masses), $\nabla_{u} \score$ blows up as $\mathcal{O}(1/\tau)$ in the interfacial layer, leading to $L_{s}(\tau) \sim 1/\tau$ and hence $\amp$ diverges (or becomes large, depending on the early-stopping $\tau_{\mathrm{end}}$). This exponential amplification is the \textbf{inherent deficiency} of standard score-based diffusion when handling discontinuous solutions.

Theorem~3 reduces this amplification factor to $\mathcal{O}(1)$ via BV-aware parameterization, constituting the core contribution of this paper.
\end{remark}

\begin{remark}[Relation to Score Shocks \citep{sarkar2026scoreshocks}]
\label{rem:score-shocks-relation}
The bound structure of this theorem is consistent with Score Shocks Theorem~6.3 ($\Wass{1} \lesssim \varepsilon \exp(\amp T)$). Our independent derivation supplements the following technical points not detailed in Score Shocks:
\begin{itemize}
\item Formal treatment of the measure mismatch (Step~5, Assumption~A4);
\item Explicit tracing of the Lipschitz constant (Assumption~A1);
\item Bridging to Kruzhkov $L^{1}$-contractivity, enabling the bound to transfer from the diffusion model domain to the PDE domain.
\end{itemize}
\end{remark}

\begin{remark}[The constant $C'$ and the network Lipschitz bound]
\label{rem:C-prime}
The auxiliary constant $C' = C_{\mathrm{stab}}\,L_{g}$ that appears in Steps~5--6 accumulates as $e^{C' T_{d}}$ in the exponent. At practical scales ($T_{d} \sim 1$, $L_{g}$ controlled by spectral normalization or weight decay), $e^{C' T_{d}} = \mathcal{O}(1)$, so this constant does not affect the asymptotic rate but only the absolute constant $C$. Formally, one may require the network to satisfy $\norm{\sth}_{\mathrm{Lip}} \le L_{\mathrm{net}}$ as part of training regularization.
\end{remark}
